\documentclass[../../main.tex]{subfiles}

\begin{document}
\begin{CJK*}{UTF8}{gbsn}

\section*{Exercise 3}

Given data set $(X_1, Y_1), \cdots, (X_{n}, Y_{n}), (X_{n+1}, Y_{n+1})
\in \mathcal{X} \times \mathcal{Y}$ with $\mathcal{Y} = \mathbb{R}$,
level $\alpha \in (0,1)$, we discretize $[-B,B]$, where $B > 0$,
by setting, for a positive integer $M$ and all $m \in \{1, \ldots, M\}$:

\begin{equation*}
    y^{(m)} = - B + \frac{m-1}{M-1} \times 2B.
\end{equation*}
Put, for any $y \in \mathcal{Y}$:

\begin{equation*}
    \mathcal{D}_{n+1}^y = \curly{(X_1, Y_1), \cdots, (X_{n}, Y_{n}), (X_{n+1}, y)}.
\end{equation*}
With a real valued score function $s((x,y); \mathcal{D})$
for $(x,y) \in \mathcal{X} \times \mathcal{Y}$
and $\mathcal{D}$ is any data set containing points from $\mathcal{X} \times \mathcal{Y}$,
we consider a discretized conformal prediction set defined as follows.
First of all, let $I$ be the set of all $m \in \{1, \ldots, M\}$
such that 

\begin{equation*}
    s \bracket{(X_{n+1}, y^{(m)}); \mathcal{D}_{n+1}^{y^{(m)}}} 
    \leqslant
    \text{Quantile} \bracket{\curly{s \bracket{(X_{i}, Y_i); \mathcal{D}_{n+1}^{y^{(m)}}}}_{i=1}^n; 
    (1- \alpha) \bracket{1 + \frac{1}{n}}}.
\end{equation*}
Then, put 

\begin{equation*}
    \mathcal{C}(X_{n+1}) = \bigcup_{m \in I} (y^{(m-1)}, y^{(m+1)}).
\end{equation*}

Construct an example of a symmetric score function $s$ and 
a distribution $P$ of the data set 
such that 

\begin{equation*}
    P \curly{Y_{n+1} \in \mathcal{C}(X_{n+1})} = 0
\end{equation*}
for any choice of $B$ and $M$.

\paragraph{Solution}
We construct the example as follows.
First of all, map the set $\mathcal{X}$ to $\mathbb{R}$ in any way one wants.
From now on, we shall identify $\mathcal{X}$ with $\mathbb{R}$
and we will write for simplicity $\mathcal{X} = \mathbb{R}$.
We construct it such that the pairs $(X_i,Y_i)$ ranged over $i \in \{1, \ldots, n+1\}$
are independent identically distributed with distribution of $X_i$'s being 
any fixed continuous probability measure (say, standard normal) and $Y_i = X_i$.
As of the score, we simply put

\begin{equation*}
    s((x,y); \mathcal{D}) = \abs{x-y}
\end{equation*}
for all $(x,y)$ and all data set $\mathcal{D}$.
Clearly, $s$ is symmetric because it does not even depend on $\mathcal{D}$.
If $i \in \{1, \ldots, n\}$, our construction yields, almost surely, that 

\begin{equation*}
    s \bracket{(X_{i}, Y_i); \mathcal{D}_{n+1}^{y^{(m)}}} = \abs{X_i-Y_i} = 0
\end{equation*}
and
\begin{equation*}
    \text{Quantile} \bracket{\curly{s \bracket{(X_{i}, Y_i); \mathcal{D}_{n+1}^{y^{(m)}}}}_{i=1}^n; 
    (1- \alpha) \bracket{1 + \frac{1}{n}}} = 0
\end{equation*}
Therefore, $I = \emptyset$ and $\mathcal{C}(X_{n+1}) = \emptyset$
for all choice of $M$ and all choice of $B$.
////

\end{CJK*}
\end{document}
